\section{Inverting the corrected linearization near the reference}
\label{sec:stable-inverse}
\label{sec:grn-fixed-phase-green}

The goal of this section is
Proposition~\ref{prop:rewrite-volume-inverse}: for every state $\bx$
in the neighbourhood of
Proposition~\ref{prop:uniform-physical-geometry}, the operator
$\widehat L_{\bx}$, bordered by the derivative of the equation with
respect to the metric parameter $\nu$, has an inverse which satisfies
tame estimates.

The argument has to respect one constraint. The reference inverse
loses coefficient weight, and a tame estimate with a loss is not a
bounded inverse on a single Banach space; it cannot by itself justify
a Neumann series. We therefore proceed in three steps.
\begin{enumerate}[label=\textup{\arabic*.}]
\item We put the problem in canonical form and identify a single
 Hilbert space of sources on which the \emph{relative} operator
 $\mathscr R_{\bx}=(\mathbb P_{\bx}-\mathbb P_*)\,\mathbb G_*$ is
 bounded, with small norm. The loss of the reference inverse is
 compensated by the factor $H^\beta$ carried by every difference of
 coefficients
 (Sections~\ref{subsec:canonical-problem}--\ref{subsec:relative}).
\item A Neumann series on that one space gives the solution.
\item All higher regularity of this same solution is then recovered
 from the equation: angular derivatives by commutators, radial
 derivatives by solving the equation for its highest radial
 derivative (Section~\ref{subsec:higher}).
\end{enumerate}
In this order of argument, one neighbourhood of the reference state
serves for all orders of differentiation.

\subsection{The linear problem in canonical form}
\label{subsec:canonical-problem}

\paragraph{Transformations.}
Use the cylinders of Section~\ref{sec:fixed-geometry}. Let
$\Psi_{\bx}$ denote the coordinate map on each piece, $J_{\bx}$ its
Jacobian determinant and $h_{\bx}$ the height factor in the original
variables. Define
\[
 \mathscr C_{\bx}u=\frac{u}{h_{\bx}}\circ\Psi_{\bx}^{-1},
 \qquad
 \mathscr S_{\bx}f=\frac{h_{\bx}f}{|J_{\bx}|}\circ\Psi_{\bx}^{-1},
 \qquad
 \mathbb L_{\bx}=
 \mathscr S_{\bx}\,\widehat L_{\bx}\,\mathscr C_{\bx}^{-1} .
\]
Thus $\mathscr C_{\bx}$ maps a height to its canonical unknown,
$\mathscr S_{\bx}$ maps a source to its canonical source, and
$\mathbb L_{\bx}$ is the operator in canonical form; on a positive
cylinder $\mathbb L_{\bx}=-P_{\bx}$. The pairing with $dr\,d\theta$ is
preserved. A \emph{piece} is one of the cylinders: a half of a
negative component, or a positive component. A function on the
collection of all pieces is written $F=(F_m)$.

\paragraph{Source norms.}
For real $t$ and integers $s\ge0$ let $\mathbb V_{t,s}$ be the space
of $F=(F_m)$ with
\begin{equation}\label{eq:vol-source-spaces}
 \|F\|_{t,s}^2
 =\sum_m M_m^{2t}\,H_m^{-2\beta}\,
                  \|F_m\|_{L^2_{\rm rad}H^s_\varphi}^2<\infty ,
\end{equation}
where $m$ runs over all pieces, the norm is $L^2$ in the radial
variable of the cylinder with $s$ angular derivatives, and $M_m,H_m$
are the weights of the cell containing the piece. No matching
condition and no integral condition is imposed on sources. The index
$t$ counts powers of the weight $M$, the index $s$ angular
derivatives.

\paragraph{Scalars and the augmented operator.}
The linear problem involves three scalars besides the canonical
unknown $W$.
\begin{itemize}
\item \emph{The endpoint constant $a_L$.} A height is represented as
 \[
  \eta=\mathscr C_{\bx}^{-1}W+a_L\chi_L,\qquad \ell_*^cW=0 ,
 \]
 so that $\eta$ tends to $a_L$ at the left endpoint and to $0$ at the
 right one. The coordinates agree with the reference ones on the
 interval $J_*$ of Section~\ref{subsec:positive-crossing}, which
 contains the support of $g_*$; hence the normalization
 $\ell_*\eta=0$ takes a form $\ell_*^cW=0$ which does not depend on
 the state.
\item \emph{The metric increment $\mu$.} It multiplies the column
 \[
  k_{\bx}=D_\nu\mathcal F_q(w,u),\qquad
  k_{\bx}^c=\mathscr S_{\bx}k_{\bx},
 \]
 where $D_\nu$ is the derivative of the residual with respect to
 $\nu$ for fixed $u$ and $q$; it acts on the metric and on the
 radial height $h$. At the reference state $k_*=g_*e_2$, the border
 column of Proposition~\ref{prop:rad-fixed-theorem}. This scalar
 removes the obstruction in the second mode.
\item \emph{An auxiliary scalar $\lambda$.} Choose a function $\psi$
 of $r$, compactly supported in $J_*$, with $\Qint\psi=1$, and put
 $c_*=\mathscr S_*\psi=\mathscr S_{\bx}\psi$. The scalar $\lambda$
 multiplies $c_*$ and removes the integral obstruction, so that
 arbitrary sources can be treated. For compatible sources it will
 turn out to be zero.
\end{itemize}
With $s_{\bx}=\mathscr S_{\bx}\widehat L_{\bx}\chi_L$, the
\emph{augmented operator} is
\begin{equation}\label{eq:vol-abstract-complete-operator}
 \mathbb P_{\bx}\bW
 =\mathbb L_{\bx}W+a_L\,s_{\bx}+\mu\,k_{\bx}^c+\lambda\,c_*,
 \qquad \bW=(W,a_L,\mu,\lambda) .
\end{equation}
The norm of an output $\bW$ is the norm \eqref{eq:vol-source-spaces}
of $W$ plus the absolute values of the three scalars, which are
assigned to the finite middle part; we write
$\|\bW\|^{\rm vol}_{t,s}$.

\paragraph{The reference inverse in this form.}
At the reference state let $\mathfrak g_*=\mathscr C_*1$ and
$\Qint_*F=\langle\mathfrak g_*,F\rangle$; this functional is
continuous by the coefficient bounds, it is the integral $\Qint$ in
canonical form, and $\Qint_*c_*=1$. Given $F$, put
$\lambda=\Qint_*F$ and apply the inverse of
Proposition~\ref{prop:rad-fixed-theorem} to the compatible source
$F-\lambda c_*$. This defines a two-sided inverse $\mathbb G_*$ of
$\mathbb P_*$. In the zero mode it is explicit: with
$f_*=\mathscr S_*^{-1}F$,
\[
 p_0(r)=\int_{r_-}^r\bigl(\av{f_*}-\lambda\psi\bigr)(t)\,dt,
 \qquad
 u_0(r)=-\int_r^{r_+}\frac{p_0(t)}{w_*(t)}\,dt ,
\]
as in \eqref{eq:rad-eq-rm49}; near either endpoint $p_0$ is the
primitive starting at that endpoint, division by $w_*$ costs a
coefficient power and not an endpoint exponent, and subtracting the
left endpoint value times $\chi_L$ gives the field and the scalar
$a_L$.

\subsection{Two endpoint exponents give decay off the diagonal}
\label{subsec:two-exponents}

By Proposition~\ref{prop:rad-fixed-theorem}(iv),(v) the reference
inverse is bounded with the endpoint weight $H^{-\gamma}$ for every
$\gamma$ in a compact subinterval of $(0,1)$. Fix
\[
 0<\epsilon<\min(\beta,1-\beta).
\]
The inverses at the exponents $\beta-\epsilon$ and $\beta+\epsilon$
coincide on smooth compactly supported sources: their difference is a
homogeneous solution in the graph domain, hence zero by
Lemma~\ref{lem:rad-green-completion}. So both estimates apply to the
same operator $\mathbb G_*$. Using each of them on one side of the
diagonal converts the two bounds into decay in the distance between
source and output (Figure~\ref{fig:two-weights}).

\begin{lemma}[A fixed loss of coefficient weight]
\label{lem:two-weights-closed-domain}
There is a fixed $L_0$ such that for every real $t$ and every
$s\ge0$
\begin{equation}\label{eq:reference-negative-material-extension}
 \|\mathbb G_*F\|^{\rm vol}_{t-L_0,s}\le C_{t,s}\,\|F\|_{t,s} .
\end{equation}
The first radial derivative of $\mathbb G_*F$ satisfies the same
bound with at most one additional angular derivative of the source.
For a fixed integer $b\ge2$,
\begin{equation}\label{eq:fixed-base-graph-domain}
 \mathbb D_*:=\mathbb G_*(\mathbb V_{0,b}),\qquad
 \|\bW\|_{\mathbb D_*}=\|\mathbb P_*\bW\|_{0,b} ,
\end{equation}
is a Hilbert space on which $\mathbb P_*$ is closed. It contains
every smooth $\bW$ in the transmission domain of
Section~\ref{subsec:transmission-domain} which satisfies the
conditions at the endpoints and the normalization, and for which
$\mathbb P_*\bW\in\mathbb V_{0,b}$.
\end{lemma}

\begin{proof}
\step{1: blocks}
Let $\bar S_m=S_m-S_N$ on the tails and $\bar S_m=0$ on the finite
middle. Let $\pi_k$ be the restriction to the piece $k$ and $\iota_m$
the extension by zero of a source on the piece $m$. By
Section~\ref{subsec:interface}, the estimate at the exponent $\gamma$
gives
\[
 \|\pi_k\,\mathbb G_*\,\iota_m\|_{H^s_\varphi\to H^s_\varphi}
 \le C_s\,\mathcal E_k^{C_s}\,\mathcal E_m^{C_s}\,
                        e^{-\gamma(\bar S_k-\bar S_m)} .
\]
There is no loss of angular derivatives, because the reference
inverse acts on each angular mode with bounds uniform in the mode.
Multiply by the weights $H_k^{-\beta}H_m^{\beta}$ of
\eqref{eq:vol-source-spaces}, and use $\gamma=\beta+\epsilon$ when
$\bar S_k\ge\bar S_m$ and $\gamma=\beta-\epsilon$ otherwise. The
resulting block $K_{km}$ satisfies
\begin{equation}\label{eq:global-two-weight-block-bound}
 \|K_{km}\|_{H^s_\varphi\to H^s_\varphi}
 \le C_s\,\mathcal E_k^{C_s}\,\mathcal E_m^{C_s}\,
                    e^{-\epsilon|\bar S_k-\bar S_m|} .
\end{equation}
The estimate for the first radial derivative in
Proposition~\ref{prop:rad-fixed-theorem}(iv) gives the same block
bound with one more angular derivative.

\step{2: summation}
For every $\delta>0$,
$|\log M_k-\log M_m|\le\delta\,|\bar S_k-\bar S_m|+C_\delta$, because
the derivative of $n^{5/2}\log^2(e+n)$ with respect to $2n^3/3$ tends
to zero. With $\mathcal E^{C_s}\le C_sM$ we get, for $L_0\ge2$,
\[
 M_k^{t-L_0}\,\|K_{km}\|\,M_m^{-t}
 \le C_s\,M_k^{2-L_0}\Bigl(\frac{M_k}{M_m}\Bigr)^{t-1}
                         e^{-\epsilon|\bar S_k-\bar S_m|}
 \le C_{t,s}\,e^{-(\epsilon/2)|\bar S_k-\bar S_m|} ,
\]
where $\delta$ was chosen with $|t-1|\,\delta\le\epsilon/2$. This
matrix has bounded row and column sums, since the numbers $S_n$ are
uniformly separated and each index carries a bounded number of
pieces. Schur's test proves
\eqref{eq:reference-negative-material-extension}.

\step{3: the domain}
By definition $\mathbb D_*$ is isometric to $\mathbb V_{0,b}$, hence
complete. If $\bW_i\to\bW$ in the space of outputs and
$\mathbb P_*\bW_i\to F$ in $\mathbb V_{0,b}$, then
$\bW_i=\mathbb G_*\mathbb P_*\bW_i\to\mathbb G_*F$, so
$\bW=\mathbb G_*F$: the operator is closed. Approximating a source by
smooth compactly supported ones, the estimate gives convergence of
the solutions on every piece, the equation passes to the limit in the
sense of distributions, and the bound for the first radial derivative
identifies the traces, so the transmission conditions hold in the
weak sense. Finally, a smooth $\bW$ as in the last statement differs
from $\mathbb G_*\mathbb P_*\bW$ by a normalized homogeneous solution
with the natural behaviour at the endpoints, which vanishes by
Lemma~\ref{lem:rad-green-completion}.
\end{proof}

\begin{figure}[htbp]
\centering
\begin{tikzpicture}[x=.92cm,y=.9cm,font=\small,>=Stealth]
\draw[->] (0,0)--(6.2,0) node[right] {depth $\bar S_m$ of the source};
\draw[->] (0,0)--(0,4.9) node[above] {depth $\bar S_k$ of the output};
\fill[blue!7] (0,0)--(4.5,4.5)--(0,4.5)--cycle;
\fill[orange!9] (0,0)--(5.7,0)--(5.7,4.5)--(4.5,4.5)--cycle;
\draw[thick] (0,0)--(4.5,4.5);
\node at (1.45,3.4) {use $\gamma=\beta+\epsilon$};
\node at (4,1.2) {use $\gamma=\beta-\epsilon$};
\node[align=left,anchor=west] at (6.55,2.6)
 {after conjugation by $H^{-\beta}$:\\[2pt]
  $e^{-\epsilon|\bar S_k-\bar S_m|}$};
\end{tikzpicture}
\caption{The same reference inverse is estimated with two endpoint
weights, each used on one side of the diagonal. Together they give
decay in the distance between the depths of the source and of the
output.}
\label{fig:two-weights}
\end{figure}

\subsection{Positive cylinders: an estimate without boundary conditions}
\label{subsec:cubic-estimate}

On a positive cylinder the perturbation
\eqref{eq:readable-positive-difference} is of second order, the same
order as the operator. It can nevertheless be controlled by the
operator, because its coefficients vanish to third order at the two
ends. The following estimate is the reason for that requirement. In
this subsection $\langle j\rangle=(1+j^2)^{1/2}$.

\begin{lemma}[Weighted graph estimate]\label{lem:cubic}
On $0<y<1$ consider, for an integer $j$,
\[
 P_ju=-(a\,u')'+j^2c\,u+V\,u=F,
 \qquad a\ge\mathcal E_m^{-C},\quad
 c\ge\mathcal E_m^{-C}\,\bd,\quad \bd=y(1-y),
\]
where the derivatives up to order two of $a,c,V$, and $a^{-1}$, are
bounded by $\mathcal E_m^{C_1}$. For every solution
$u\in L^2$ with $F\in L^2$, without any condition at the endpoints,
\begin{equation}\label{eq:vol-arbitrary-boundary-graph-bound}
 \begin{split}
 &\|\bd\,u'\|_2+\langle j\rangle\,\|\bd^{3/2}u\|_2+\|\bd^{2}u''\|_2
   +\langle j\rangle^2\,\|\bd^{3}u\|_2
   +\langle j\rangle\,\|\bd^{5/2}u'\|_2\\
 &\hspace{30mm}\le\mathcal E_m^C\,\bigl(\|F\|_2+\|u\|_2\bigr).
 \end{split}
\end{equation}
\end{lemma}

\begin{proof}
For fixed $j$ a solution in $L^2$ with $F\in L^2$ is in $H^2$ up to
the endpoints, since the equation is a regular ordinary differential
equation; this justifies the integrations by parts below, but the
uniform bounds come from the computation.

\step{1: first order}
Multiply the equation by $\bd^2\bar u$ and integrate. Since $\bd$
vanishes at both ends,
\[
 \int a\,|(\bd u)'|^2+j^2\int c\,\bd^2|u|^2
 =\operatorname{Re}\langle F,\bd^2u\rangle
       +\int a\,|\bd'|^2|u|^2-\int V\,\bd^2|u|^2 .
\]
This bounds $\bd\,u'$ and $j\,\bd^{3/2}u$.

\step{2: second order}
The function $\widetilde u=\bd^2u$ vanishes at both ends and satisfies
\[
 P_j\widetilde u=\bd^2F-2a\,(\bd^2)'u'-(a\,(\bd^2)')'u ,
\]
whose right side is controlled by Step~1. With $\mathsf c=c/a$ the
same is true of
$F_1=-\widetilde u''+j^2\mathsf c\,\widetilde u$. Squaring and
integrating the cross term by parts twice,
\[
 \|\widetilde u''\|_2^2+j^4\|\mathsf c\,\widetilde u\|_2^2
 +2j^2\int\mathsf c\,|\widetilde u'|^2
 =\|F_1\|_2^2+j^2\int\mathsf c''\,|\widetilde u|^2 .
\]
The last integral is controlled by Step~1 because
$\bd^4\le C\bd^3$. Using $\mathsf c\ge\mathcal E_m^{-C}\bd$ and
expressing $\widetilde u'$ and $\widetilde u''$ through $u$ gives
\eqref{eq:vol-arbitrary-boundary-graph-bound}.
\end{proof}

\paragraph{Consequence for the perturbation.}
Every second-order term of \eqref{eq:readable-positive-difference} is
controlled by the left side of
\eqref{eq:vol-arbitrary-boundary-graph-bound}: $\bd^3u_{yy}$ by
$\bd^2u''$, $\bd^3u_{y\varphi}$ by
$\langle j\rangle\bd^{5/2}u'$, and $\bd^3u_{\varphi\varphi}$ by
$\langle j\rangle^2\bd^3u$. When a derivative falls on a coefficient,
the remaining terms are controlled by the first two norms.

The coefficients depend on $\varphi$, and this dependence is handled
by Fourier series. Write $u=\sum_ju_j(y)\,e^{ij\varphi}$ and, for one
of the coefficients, $e=\sum_ke_k(y)\,e^{ik\varphi}$. The $l$th
Fourier coefficient of $\bd^3e\,u_{yy}$ is
$\sum_ke_k\,\bd^3u_{l-k}''$, so
\[
 \bigl\|(\bd^3e\,u_{yy})_l\bigr\|_{L^2_y}\le\sum_k\epsilon_k\,g_{l-k},
 \qquad \epsilon_k=\sup_y|e_k(y)|,\quad
 g_j=\|\bd^2u_j''\|_{L^2_y}
 \le\mathcal E_m^C\bigl(\|F_j\|_2+\|u_j\|_2\bigr),
\]
the last inequality being Lemma~\ref{lem:cubic} for the mode $j$.
Since $\langle l\rangle^s\le2^s(\langle k\rangle^s+\langle l-k\rangle^s)$,
Young's inequality for sequences gives
\begin{equation}\label{eq:sequence-convolution}
 \bigl\|\langle l\rangle^s(\epsilon*g)_l\bigr\|_{\ell^2}
 \le2^s\bigl(\|\epsilon\|_{\ell^1}\,\|\langle j\rangle^sg\|_{\ell^2}
 +\|\langle k\rangle^s\epsilon\|_{\ell^2}\,\|g\|_{\ell^1}\bigr),
\end{equation}
and $\|g\|_{\ell^1}\le C\|\langle j\rangle g\|_{\ell^2}$,
$\|\epsilon\|_{\ell^1}\le C\|\langle k\rangle\epsilon\|_{\ell^2}$.
The weighted norms of $(\epsilon_k)$ are bounded by norms of the
coefficient $e$ with one more derivative in $y$, hence by
\eqref{eq:geometry-current-jet-bound}. The other terms are treated in
the same way, with $g_j$ replaced by the corresponding term on the
left of \eqref{eq:vol-arbitrary-boundary-graph-bound}; for instance
$g_j=\langle j\rangle^2\|\bd^3u_j\|_{L^2_y}$ for the term
$\bd^3e\,u_{\varphi\varphi}$. For the commutator with
$\Lambda=|D_\varphi|$, the $l$th coefficient of $[\Lambda,e]\,v$ is
$\sum_k(|l|-|l-k|)\,e_kv_{l-k}$, and
$\bigl||l|-|l-k|\bigr|\le|k|$; so $\epsilon_k$ is replaced by
$|k|\,\epsilon_k$, which costs one more angular derivative of the
coefficient. Altogether, for $P_*u=F$ on the piece $m$ and $s\ge b$,
\begin{equation}\label{eq:readable-positive-relative-bound}
 \|\delta P_{\bx}u\|_s
 \le C_s\,H_m^\beta\,\mathcal E_m^{C_s}
 \Bigl[\mathfrak g_{b_0,m}\,\bigl(\|F\|_s+\|u\|_s\bigr)
       +\mathfrak g_{s+d_*,m}\,\bigl(\|F\|_b+\|u\|_b\bigr)\Bigr],
\end{equation}
and the same bound holds for $[\Lambda,\delta P_{\bx}]u$. Here
$\delta P_{\bx}=P_{\bx}-P_*$, the norms are those of
$L^2_yH^s_\varphi$, and $\mathfrak g_{s,m}$ is the localized seminorm
of Section~\ref{subsec:neighbourhood}. No artificial boundary
condition was imposed on the reference solution $u$.

\paragraph{An elementary estimate for the first radial derivative.}
On every piece, normalized to unit radial length, the equation at
frequency $j$ has the form $Au''+Bu'+(j^2C+D_0)u=F$, with $A$ bounded
away from zero and all coefficients bounded by coefficient powers.
From the equation and the interpolation inequality
$\|u'\|_2^2\le C\|u\|_2(\|u''\|_2+\|u\|_2)$,
\begin{equation}\label{eq:vol-reference-first-normal-bound}
 \|u'\|_2\le\mathcal E_m^C\bigl(\|F\|_2+\langle j\rangle\|u\|_2\bigr),
 \qquad
 \|u''\|_2\le\mathcal E_m^C\bigl(\|F\|_2+\langle j\rangle^2\|u\|_2\bigr),
\end{equation}
with no assumption on the boundary values: substitute the equation
for $u''$ into the interpolation inequality and absorb the term
$\mathcal E_m^C\|u\|_2\|u'\|_2$ by Young's inequality. Smooth sources
with finitely many frequencies are dense, and for limits of such
sources these estimates identify the weak derivatives; so no
regularity is assumed beforehand.

\subsection{The relative operator}
\label{subsec:relative}

On smooth compactly supported sources define
\[
 \mathscr R_{\bx}=(\mathbb P_{\bx}-\mathbb P_*)\,\mathbb G_* .
\]
If $\mathbb G_*F=(W,a_L,\mu,\lambda)$, then
\begin{equation}\label{eq:vol-normalized-relative-identity}
 \mathscr R_{\bx}F
 =\delta\mathbb L_{\bx}W
      +a_L\,(s_{\bx}-s_*)+\mu\,(k_{\bx}^c-k_*^c),
 \qquad\delta\mathbb L_{\bx}=\mathbb L_{\bx}-\mathbb L_* ;
\end{equation}
the column $c_*$ does not depend on the state. On a negative piece
$\delta\mathbb L_{\bx}$ is the difference of the potentials $V_i$ in
\eqref{eq:fixed-strip-airy-equation}, a multiplication operator. On a
positive piece it is estimated by
\eqref{eq:readable-positive-relative-bound}, the compactly supported
scalar columns being counted as part of the source. In every case the
difference of coefficients on the output piece $k$ carries the factor
$H_k^\beta$ of \eqref{eq:geometry-current-jet-bound}.

\begin{lemma}[The relative operator is small on one space]
\label{lem:relative}
Let $b\ge2$ be as in Lemma~\ref{lem:two-weights-closed-domain}. After
$b_0$ is increased and $\epsilon_0$ decreased,
\begin{equation}\label{eq:rewrite-base-relative}
 \|\mathscr R_{\bx}\|_{\mathbb V_{0,b}\to\mathbb V_{0,b}}
                         \le C\,\mathfrak n_{b_0}<\tfrac12 .
\end{equation}
Moreover, for $t\ge0$ and $s\ge b$, with $\Lambda=|D_\varphi|$,
\begin{equation}\label{eq:compact-volume-relative-bound}
 \|\mathscr R_{\bx}F\|_{t,s}
 +\|[\Lambda,\mathscr R_{\bx}]F\|_{t,s}
 \le C_{t,s}\bigl[\epsilon_0\,\|F\|_{0,s}
             +\mathfrak n_{t+s+d_*}\,\|F\|_{0,b}\bigr] .
\end{equation}
The operator $\mathscr R_{\bx}$ extends by continuity to
$\mathbb V_{0,b}$, and
\begin{equation}\label{eq:current-fixed-domain-factorization}
 \mathbb P_{\bx}=(I+\mathscr R_{\bx})\,\mathbb P_*
                         \qquad\text{on }\mathbb D_* .
\end{equation}
\end{lemma}

Note that \eqref{eq:rewrite-base-relative} is the boundedness of an
operator from one space to the same space. No operator that loses
derivatives or weights is treated as a bounded endomorphism.

\begin{proof}
\step{1: the base space}
Combine the factor $H_k^\beta$ with
\eqref{eq:global-two-weight-block-bound}. With coefficients measured
at low order, the block of $\mathscr R_{\bx}$ from the piece $m$ to
the piece $k$ has norm at most
\begin{equation}\label{eq:rewrite-relative-block}
 C\,\mathfrak n_{b_0}\,M_k^C\,M_m^C\,
            e^{-\beta\bar S_k-\epsilon|\bar S_k-\bar S_m|} .
\end{equation}
For every fixed $C$ and $\delta>0$ we have
$M_i^C\le C_\delta\,e^{\delta\bar S_i}$, and
$\bar S_m\le\bar S_k+|\bar S_k-\bar S_m|$. Choosing
$\delta<\min(\beta/4,\epsilon/2)$ bounds
\eqref{eq:rewrite-relative-block} by
\[
 C'\,\mathfrak n_{b_0}\,
       e^{-(\beta/2)\bar S_k-(\epsilon/2)|\bar S_k-\bar S_m|} ,
\]
which has bounded row and column sums. Increasing $b_0$ to cover the
finitely many low derivatives used, and then decreasing $\epsilon_0$,
gives \eqref{eq:rewrite-base-relative}.

\step{2: higher orders}
In the product estimates keep one coefficient factor at the high
order. The same summation proves the bound for $\mathscr R_{\bx}$ in
\eqref{eq:compact-volume-relative-bound}; the output may carry any
power $M^t$ because of the factor $H_k^\beta$. For the commutator,
assign the frequency zero to the scalars $a_L,\lambda$ and the
frequency two to $\mu$. For multiplication by a coefficient, the
matrix of the commutator in Fourier modes gains the factor
$|k|-|j|$, which is bounded by $|k-j|$; so one more derivative of the
coefficient gives the same bound. The commutators with the scalar
columns are the terms $\Lambda(s_{\bx}-s_*)\,a_L$ and
$(\Lambda-2)(k_{\bx}^c-k_*^c)\,\mu$, which are of the same type.

\step{3: the factorization}
By Step~1 the operator extends to $\mathbb V_{0,b}$. Approximating by
smooth compactly supported sources identifies the extension with the
difference of the differential operators, in the sense of
distributions, on $\mathbb D_*$; here it is used that the
transmission conditions are the same for all states. This is
\eqref{eq:current-fixed-domain-factorization}.
\end{proof}

\subsection{Solving, and recovering all derivatives}
\label{subsec:higher}

By \eqref{eq:rewrite-base-relative} we may define, for
$F\in\mathbb V_{0,b}$,
\[
 \widetilde F=(I+\mathscr R_{\bx})^{-1}F,
 \qquad \bW=\mathbb G_*\widetilde F ,
\]
where the inverse is a Neumann series on $\mathbb V_{0,b}$. By
\eqref{eq:current-fixed-domain-factorization},
$\mathbb P_{\bx}\bW=F$. We now show that this solution is as regular
as the source allows, with tame estimates. Three steps recover, in
turn, the angular derivatives of $\widetilde F$, its coefficient
weight, and the radial derivatives of $W$.

\paragraph{Step A: angular derivatives of $\widetilde F$.}
This step rests on a general lemma. Let $\mathcal H$ be a Hilbert
space with a self-adjoint operator $\Lambda$ whose spectrum lies in
$\{0,1,2,\ldots\}$, let $\|v\|_p=\|\langle\Lambda\rangle^pv\|$, and
let $\mathcal H^p$ be the space of vectors for which this norm is
finite.

\begin{lemma}[Regularity of an inverse from a commutator]
\label{lem:ang-base-inverse-regularity}
Let $R$ be a bounded operator on $\mathcal H$ which maps every
$\mathcal H^p$ into itself, and suppose that $I+R$ has a bounded
inverse on $\mathcal H$. Suppose the commutator
$\mathscr K=[\Lambda,R]$ satisfies, for all integers $p\ge0$,
\begin{equation}\label{eq:commutator-hypothesis}
 \|\mathscr Kv\|_p\le a_p\,\|v\|_p+c_p\,\|v\|_0 ,
\end{equation}
with nondecreasing constants $a_p$ and $c_p$. Then $(I+R)^{-1}$ maps every
$\mathcal H^p$ into itself, and for $p\ge1$
\begin{equation}\label{eq:commutator-conclusion}
 \|(I+R)^{-1}f\|_p\le C_p\bigl(\|f\|_p+(1+c_{p-1})\,\|f\|_0\bigr),
\end{equation}
where $C_p$ depends only on $p$, on $a_0,\ldots,a_{p-1}$, on $c_0$
and on the norm of $(I+R)^{-1}$ on $\mathcal H$.
\end{lemma}

The constants $c_p$ with $p\ge1$ may be arbitrarily large: they enter
linearly, and they do not enter the constant $C_p$.

\begin{proof}
\step{1: the solution is regular}
Let $U_\delta=e^{i\delta\Lambda}$. For $v\in\mathcal H^1$ the function
$t\mapsto U_{\delta-t}RU_tv$ is differentiable with derivative
$-i\,U_{\delta-t}\mathscr KU_tv$, so
\[
 [U_\delta,R]=i\int_0^\delta U_{\delta-t}\,\mathscr K\,U_t\,dt ,
\]
first on $\mathcal H^1$ and then on $\mathcal H$, because
$\mathscr K$ is bounded on $\mathcal H$ by
\eqref{eq:commutator-hypothesis} with $p=0$. Let $u=(I+R)^{-1}f$.
Then
\[
 (I+R)\,\frac{U_\delta u-u}{\delta}
 =\frac{U_\delta f-f}{\delta}-\frac{[U_\delta,R]\,u}{\delta}.
\]
If $f\in\mathcal H^1$ the right side is bounded in $\mathcal H$
uniformly in $\delta$, hence so are the difference quotients of $u$;
by Fatou's lemma for the spectral decomposition of $\Lambda$,
$u\in\mathcal H^1$, and letting $\delta\to0$ gives
$(I+R)\Lambda u=\Lambda f-\mathscr Ku$. Now let $p\ge2$, $f\in\mathcal H^p$
and suppose $u\in\mathcal H^{p-1}$. Then
\[
 (I+R)\,\Lambda^{p-1}u=\Lambda^{p-1}f
 -\sum_{k=0}^{p-2}\Lambda^k\,\mathscr K\,\Lambda^{p-2-k}u ,
\]
and the right side lies in $\mathcal H^1$, because
$\Lambda^{p-2-k}u\in\mathcal H^{k+1}$ and $\mathscr K$ maps
$\mathcal H^{k+1}$ into itself. By the case just treated
$\Lambda^{p-1}u\in\mathcal H^1$, that is, $u\in\mathcal H^p$. By
induction, for $f\in\mathcal H^p$,
\begin{equation}\label{eq:commuted-equation}
 (I+R)\,\Lambda^pu=\Lambda^pf
 -\sum_{k=0}^{p-1}\Lambda^k\,\mathscr K\,\Lambda^{p-1-k}u .
\end{equation}

\step{2: the estimate}
For $p=1$, \eqref{eq:commuted-equation} and the bound $a_0+c_0$ for
$\mathscr K$ on $\mathcal H$ give the claim. Let $p\ge2$ and
$A=1+c_{p-1}$. By \eqref{eq:commutator-hypothesis}, $\mathscr K$ is
bounded on $\mathcal H$ with norm at most $a_0+c_0$, and from
$\mathcal H^{p-1}$ with the norm
$(\|v\|_{p-1}^2+A^2\|v\|_0^2)^{1/2}$ to $\mathcal H^{p-1}$ with norm
at most $a_{p-1}+1$. Both pairs of norms are weighted $\ell^2$ norms over
the spectrum of $\Lambda$. Interpolating between them, by the
three-lines lemma applied to
$z\mapsto\langle\Lambda\rangle^{(p-1)z}\,\mathscr K\,
(\langle\Lambda\rangle^{2(p-1)}+A^2)^{-z/2}$, and using
$(\langle j\rangle^{2(p-1)}+A^2)^{\vartheta}
\le\langle j\rangle^{2(p-1)\vartheta}+A^{2\vartheta}$ for
$0\le\vartheta\le1$, we get
\[
 \|\mathscr Kv\|_k\le C\bigl(\|v\|_k+A^{k/(p-1)}\,\|v\|_0\bigr),
 \qquad 0\le k\le p-1 .
\]
Apply this to $v=\Lambda^{p-1-k}u$. Since
$\|u\|_{p-1-k}\le\|u\|_{p-1}^{1-k/(p-1)}\|u\|_0^{k/(p-1)}$, Young's
inequality gives
\[
 A^{k/(p-1)}\,\|u\|_{p-1-k}
 \le(A\|u\|_0)^{k/(p-1)}\,\|u\|_{p-1}^{1-k/(p-1)}
 \le\|u\|_{p-1}+A\,\|u\|_0 .
\]
So every term of the sum in \eqref{eq:commuted-equation} is bounded
in $\mathcal H$ by $C(\|u\|_{p-1}+A\|u\|_0)$. Applying $(I+R)^{-1}$
to \eqref{eq:commuted-equation},
\[
 \|u\|_p\le C\bigl(\|f\|_p+\|u\|_{p-1}+A\,\|u\|_0\bigr).
\]
Finally $\|u\|_{p-1}\le\varepsilon\|u\|_p+C_\varepsilon\|u\|_0$, and
$\|u\|_0\le C\|f\|_0$; absorbing the term $\varepsilon\|u\|_p$ proves
\eqref{eq:commutator-conclusion}.
\end{proof}

We apply the lemma with $\mathcal H=\mathbb V_{0,b}$,
$\mathcal H^p=\mathbb V_{0,b+p}$, $\Lambda=|D_\varphi|$ and
$R=\mathscr R_{\bx}$. By \eqref{eq:compact-volume-relative-bound}
with $t=0$, the operator $\mathscr R_{\bx}$ maps every
$\mathbb V_{0,s}$ into itself, and its commutator with $\Lambda$
satisfies \eqref{eq:commutator-hypothesis} with fixed constants $a_p$
and with $c_p=C_p\,\mathfrak n_{b+p+d_*}$; in particular $c_0$ is
bounded on the neighbourhood. Hence
$\widetilde F=(I+\mathscr R_{\bx})^{-1}F$ has all the angular
regularity of $F$, and
\begin{equation}\label{eq:angular-resolvent-bound}
 \|\widetilde F\|_{0,s}
 \le C_s\bigl[\|F\|_{0,s}+(1+\mathfrak n_{s+d_*})\,\|F\|_{0,b}\bigr],
 \qquad s\ge b .
\end{equation}
Only one factor carries the high order, and no smallness of a high
norm of the state has been used.

\paragraph{Step B: coefficient weight of $\widetilde F$.}
Use the identity
$\widetilde F=F-\mathscr R_{\bx}\widetilde F$ for the function
already constructed. By \eqref{eq:compact-volume-relative-bound} the
last term has every coefficient weight once the angular regularity of
$\widetilde F$ is known. Hence
\begin{equation}\label{eq:full-source-tame-resolvent}
 \|\widetilde F\|_{t,s}
 \le C_{t,s}\bigl[\|F\|_{t,s}
       +(1+\mathfrak n_{t+s+d_*})\,\|F\|_{0,b}\bigr] .
\end{equation}
No second inversion is needed.

\paragraph{Step C: radial derivatives of $W$.}
The corrected source $\widetilde F$ is controlled only in its angular
derivatives. For the radial derivatives of $W$ we therefore use the
equation $\mathbb P_{\bx}\bW=F$ with the \emph{original} source.
Subtract the three scalar columns,
\[
 F^\circ=F-a_L\,s_{\bx}-\mu\,k_{\bx}^c-\lambda\,c_* ,
\]
and rescale every half of a negative component to unit radial length
with a variable $y$; the positive pieces are already normalized, and
the rescaling factors are coefficient powers. Solving the equation
for its highest radial derivative gives on every piece
\begin{equation}\label{eq:rewrite-normal-recursion-base}
 \partial_y^2W=\sum_{(i,j)\in I}A_{ij}\,\partial_y^i
                      \partial_\varphi^jW+B\,F^\circ,\qquad
 I=\{(1,1),(0,2),(1,0),(0,1),(0,0)\} .
\end{equation}
On a negative half this is the equation
$W_{\xi\xi}=\xi W_{\varphi\varphi}-VW+F^\circ$. On a positive piece
the coefficient of $\partial_y^2$ is bounded below by
$\mathcal E_m^{-C}$ on the whole neighbourhood. The coefficients
$A_{ij}$ and $B$ obey the calculus of
Proposition~\ref{prop:uniform-physical-geometry}, with a fixed loss
$d_{\rm geo}$ of derivatives of the state. Differentiating, for
$a\ge2$,
\begin{equation}\label{eq:rewrite-mixed-normal-recursion}
 \begin{split}
 \partial_y^a\partial_\varphi^b W
 ={}&\sum_{(i,j)\in I}\sum_{p=0}^{a-2}\sum_{q=0}^{b}
       \binom{a-2}{p}\binom bq\,
       (\partial_y^p\partial_\varphi^q A_{ij})\;
       \partial_y^{a-2-p+i}\partial_\varphi^{b-q+j}W\\
 &+\partial_y^{a-2}\partial_\varphi^b(B\,F^\circ) .
 \end{split}
\end{equation}
Every term on the right contains at most $a-1$ radial derivatives of
$W$ and has total order at most $a+b$, and a derivative that falls on
a coefficient lowers the order of the remaining factor by the same
amount. Iterating, the expansion terminates in angular derivatives of
$W$, angular derivatives of $\partial_yW$, and derivatives of
$F^\circ$; at total order $s$ their orders are at most $s$, $s-1$ and
$s-2$. Starting from the angular regularity and the first radial
derivative, which are known, this identity proves the existence of
all mixed derivatives in turn.

To make this quantitative, fix a base order $b_1\ge4$ and for
$s\ge b_1+2$ put
\[
 \mathfrak a_{s,m}
 =\|W_m\|_{L_y^2H_\varphi^s}
   +\|\partial_yW_m\|_{L_y^2H_\varphi^{s-1}}
   +\|F_m^\circ\|_{H^{s-2}_{y,\varphi}} .
\]
The finite expansion of
\eqref{eq:rewrite-mixed-normal-recursion} gives
\begin{equation}\label{eq:rewrite-local-mixed-bound}
 \|W_m\|_{H^s_{y,\varphi}}
 \le C_s\,\mathcal E_m^{C_s}\,
       \bigl[\mathfrak a_{s,m}
       +(1+\mathfrak n_{s+d_{\rm geo},m})\,\mathfrak a_{b_1,m}\bigr] .
\end{equation}
Indeed, in each term of the expansion the orders of the derivatives
of the coefficients and of the remaining factor add up to at most
$s$. This is the situation of Lemma~\ref{lem:one-high-factor}(ii),
applied after extension to the full cylinder: one factor is taken in
the norm of order $s$, or $s+d_{\rm geo}$ for a coefficient, and all
the others in norms of the fixed order $b_1$. Where a factor is known
only in the mixed norms of $\mathfrak a_{s,m}$, interpolate its
intermediate norms between the corresponding entries of
$\mathfrak a_{b_1,m}$ and $\mathfrak a_{s,m}$ by
Lemma~\ref{lem:one-high-factor}(i). The Sobolev embedding is used at the fixed order $b_1$
only, and reciprocals of the leading coefficient are treated in the
same way thanks to its lower bound. No norm of the state of order $s$
has to be small.

The exponent $C_s$ of the reference coefficients in
\eqref{eq:rewrite-local-mixed-bound} may grow with $s$. It is
important to absorb the \emph{complete} product only after the
expansion: $\mathcal E_m^{C_s}\le C_sM_m$ costs one power of the
weight, whereas paying for each step of the recursion separately
would not give a fixed loss. Define the mixed norms
\[
 \bigl(\|Z\|_{t,s}^{\rm mix}\bigr)^2
 =\sum_m M_m^{2t}\,H_m^{-2\beta}
       \sum_{a+b\le s}\|\partial_y^a\partial_\varphi^b Z_m\|_2^2 .
\]
The first entry of $\mathfrak a_{s,m}$ is controlled by
\eqref{eq:reference-negative-material-extension}. For the second,
use the bound for the first radial derivative in
Lemma~\ref{lem:two-weights-closed-domain} at the angular order $s-1$;
its one additional angular derivative requires exactly the order $s$
of $\widetilde F$. Summing \eqref{eq:rewrite-local-mixed-bound} over
the pieces yields
\begin{equation}\label{eq:rewrite-global-mixed-bound}
 \|W\|_{t,s}^{\rm mix}
 \le C_{t,s}\bigl[
   \|\widetilde F\|_{t+L_0+1,s}
   +\|F\|_{t+1,s-2}^{\rm mix}
 +(1+\mathfrak n_{t+s+d_{\rm tot}})\,\|F\|_{0,b_1}^{\rm mix}\bigr] ,
\end{equation}
where the fixed number $d_{\rm tot}$ covers $L_0+1$, $d_{\rm geo}$
and the base and embedding orders. The scalar columns are compactly
supported and satisfy the same estimates, and the scalars themselves
are bounded by the base inverse. In the term with the high norm of
the state, the weight $M^{t+L_0+1}$ is placed on the localized norm
of the state, and the low norm of the field is taken in the output
space of the reference inverse, with weight $-L_0$. The order $b_0$
of Proposition~\ref{prop:uniform-physical-geometry} is taken large
enough to control the finitely many base norms that occur; the
finitely many orders below $b_1+2$ follow from the same computation.

\paragraph{Conclusion.}
Insert \eqref{eq:full-source-tame-resolvent} into the first term of
\eqref{eq:rewrite-global-mixed-bound}. The indices of weight and of
derivatives in that term add up to $t+s+L_0+1$, and in the source
term to $t+s-1$. So an estimate for all $t+s\le r$ loses a fixed
number of orders of the grading, independent of $r$. The
transformations between canonical and original variables have fixed
tame losses by Proposition~\ref{prop:uniform-physical-geometry}, and
the jets match across the roots as explained in
Section~\ref{subsec:transmission-domain}. After increasing $d_*$ once
to include all these losses we obtain
\eqref{eq:readable-physical-inverse-tame} below.

Undoing the transformations, $\bW$ corresponds to a height $\eta$ and
scalars $\mu,\lambda$ with
\[
 \widehat L_{\bx}\eta+\mu\,k_{\bx}+\lambda\,\psi=f,
 \qquad \ell_*\eta=0 .
\]
Integrate this equation over the annulus. The fluxes of $\eta$ vanish
at the endpoints, so $\Qint\widehat L_{\bx}\eta=0$; and
$\Qint k_{\bx}=0$, by differentiating the conservation law of
Lemma~\ref{lem:scalar-integral} with respect to $\nu$. Hence
$\lambda=\Qint f$, which is zero for $f\in\mathcal Y_\beta$.

\begin{proposition}[The inverse at a nearby state]
\label{prop:rewrite-volume-inverse}
There are $b_0$, $\epsilon_0$ and $d_*$ such that for every state
with $\mathfrak n_{b_0}<\epsilon_0$ and every $f\in\mathcal Y_\beta$
there is a unique pair $G_{\bx}f=(\eta,\mu)$ with
\[
 \widehat L_{\bx}\eta+\mu\,k_{\bx}=f,\qquad
 \eta\in\mathcal X_\beta\cap\ker\ell_*,\qquad \mu\in\R ,
\]
where $\eta$ tends to zero at the right endpoint and to a constant at
the left endpoint. For every $s$, with a fixed base order $b$,
\begin{equation}\label{eq:readable-physical-inverse-tame}
 \|G_{\bx}f\|_s
 \le C_s\bigl[\|f\|_{s+d_*}
       +(1+\mathfrak n_{s+d_*})\,\|f\|_{b+d_*}\bigr] .
\end{equation}
The map $(\bx,q,f)\mapsto G_{\bx}f$ is smooth tame.
\end{proposition}

\begin{proof}
Existence and the estimate were proved above. The factorization
\eqref{eq:current-fixed-domain-factorization} on $\mathbb D_*$ shows
that $G_{\bx}$ is a two-sided inverse, and uniqueness among smooth
solutions follows because $\mathbb D_*$ contains all of them
(Lemma~\ref{lem:two-weights-closed-domain}). Continuity in the state
follows from the resolvent identity, first in the base space and
then, by the higher estimates, in every order. The coefficients form
a smooth tame family. A smooth tame family of linear maps which has
unique solutions and a continuous tame inverse has a smooth tame
inverse (Theorem~\ref{thm:hamilton-inverse} below); its hypotheses,
uniqueness and continuity, are exactly what has just been verified.
\end{proof}
